\documentclass[12pt,a4paper]{article}
\usepackage[utf8]{inputenc}
\usepackage[T1]{fontenc}
\usepackage[turkish]{babel}
\usepackage[margin=2.5cm]{geometry}
\usepackage{graphicx,float,booktabs,tabularx}
\usepackage{amsmath}
\usepackage{enumitem}
\usepackage{hyperref}
\usepackage{setspace}
\usepackage{parskip}
\usepackage{fancyhdr}

\pagestyle{fancy}
\fancyhf{}
\fancyhead[L]{\small IoT Saldırı Tespiti --- Büyük Veri}
\fancyhead[R]{\thepage}
\renewcommand{\headrulewidth}{0.4pt}
\onehalfspacing

\begin{document}

% ── KAPAK ──────────────────────────────────────────────────────────────
\begin{titlepage}
\centering
\vspace*{3cm}
{\LARGE\bfseries IoT Ağ Trafiğinde\\[6pt]Saldırı Tipi Sınıflandırması}\\[1.2cm]
{\Large Büyük Veri Projesi Raporu}\\[2.5cm]
\begin{tabular}{rl}
\textbf{Veri Seti}    & Edge-IIoTset Cyber Security Dataset \\[4pt]
\textbf{Teknolojiler} & Kafka, Spark, Delta Lake, MLflow, Streamlit \\[4pt]
\textbf{Tarih}        & Mayıs 2025 \\
\end{tabular}
\vfill
\end{titlepage}

\tableofcontents
\newpage

%=========================================================================
\section{Giriş}

Bu proje, \textbf{Edge-IIoTset} veri setini kullanarak IoT ağ trafiği üzerinde çok sınıflı (multi-class) saldırı tipi sınıflandırması yapan uçtan uca bir büyük veri pipeline'ı sunar. Kafka ile gerçek zamanlı akış, Spark ile dağıtık işleme, Delta Lake ile katmanlı depolama, 5 ML modeli, MLflow ile deney takibi ve Streamlit dashboard bileşenleri Docker Compose ile orkestre edilir.

\textbf{Amaç:} Gelen her ağ paketinin saldırı türünü (Normal, DDoS, Port Scanning, MITM, Ransomware vb.) yüksek doğrulukla sınıflandırmak.

%=========================================================================
\section{Veri Seti}

\begin{table}[H]\centering
\begin{tabular}{ll}
\toprule
Kaynak & Kaggle (Mohamed Amine Ferrag) \\
Dosya  & \texttt{ML-EdgeIIoT-dataset.csv} ($\sim$82\,MB) \\
Kolon  & 63 (TCP, UDP, HTTP, DNS, MQTT, Modbus vb.) \\
Sınıf  & 14+ saldırı tipi + Normal \\
Hedef  & \texttt{Attack\_type} \\
\bottomrule
\end{tabular}
\end{table}

Her satır bir ağ paketini temsil eder. Paket yalnızca bir protokol kullandığından, diğer protokol kolonları doğal olarak 0'dır (eksik değer değil, aktivite yok).

%=========================================================================
\section{Sistem Mimarisi}

Tüm bileşenler \texttt{docker compose up -d} ile ayağa kalkar.

\begin{table}[H]\centering
\begin{tabular}{lll}
\toprule
\textbf{Bileşen} & \textbf{Teknoloji} & \textbf{Port} \\
\midrule
Mesaj Kuyruğu   & Kafka 3.4 + Zookeeper & 9092 \\
Dağıtık İşleme  & Apache Spark 3.4.2    & 8080 \\
Depolama         & Delta Lake 2.4.0      & --- \\
ML Takip         & MLflow                & 5000 \\
Dashboard        & Streamlit             & 8501 \\
Notebook         & JupyterLab            & 8888 \\
\bottomrule
\end{tabular}
\end{table}

Spark cluster: 1 master + 1 worker (2\,GB RAM, 2 core).

%=========================================================================
\section{Veri Pipeline'ı}

\subsection{Kafka Producer}
CSV kayıtlarını JSON'a çevirip \texttt{iot-network-traffic} topic'ine (3 partition) gönderir. Her mesaja \texttt{timestamp}, \texttt{flow\_id}, IP adresleri eklenir.

\subsection{Medallion Architecture (Delta Lake)}

\begin{enumerate}[noitemsep]
  \item \textbf{Bronze:} Ham JSON'u değiştirmeden Delta'ya yazar (yıl/ay/gün partition).
  \item \textbf{Silver:} JSON parse, kolon rename, 8 gereksiz kolon drop, inf temizleme, eksik değer doldurma, tip cast, duplikat kaldırma.
  \item \textbf{Gold:} Silver + 5 yeni feature $\Rightarrow$ ML-ready tablo.
\end{enumerate}

Her katman 30\,s micro-batch ile Structured Streaming olarak çalışır.

%=========================================================================
\section{Feature Engineering}

Silver verisine eklenen 5 özellik:

\begin{table}[H]\centering
\begin{tabularx}{\textwidth}{clXl}
\toprule
\# & Feature & Formül & Hedef Saldırı \\
\midrule
1 & \texttt{traffic\_asymmetry\_ratio} & $\frac{\text{tcp\_ack}}{\text{tcp\_seq}+1}$ & DDoS Flood \\[4pt]
2 & \texttt{pkt\_size\_cv} & $\frac{\text{tcp\_len}}{|\text{seq}-\text{ack}|+1}$ & Port Scan \\[4pt]
3 & \texttt{flow\_intensity} & $\text{tcp\_len} \times \text{tcp\_flags}$ & DDoS UDP \\[4pt]
4 & \texttt{iat\_regularity} & $\frac{\text{tcp\_checksum}}{\text{tcp\_len}+1}$ & Brute Force \\[4pt]
5 & \texttt{conn\_efficiency} & $\frac{\text{syn}}{\text{fin}+\text{rst}+1}$ & SYN Flood \\
\bottomrule
\end{tabularx}
\end{table}

%=========================================================================
\section{Makine Öğrenmesi Modelleri}

\subsection{Ortak Pipeline}
\begin{enumerate}[noitemsep]
  \item Gold Delta'dan batch okuma
  \item \texttt{StringIndexer} ile Attack\_type $\to$ sayısal label
  \item \texttt{VectorAssembler} ile feature vektörü
  \item Balanced class weight: $w_c = \frac{N}{K \cdot N_c}$
  \item Stratified train/test split (\%80/\%20, seed=42)
\end{enumerate}

\subsection{Eğitilen Modeller}

\begin{table}[H]\centering
\begin{tabularx}{\textwidth}{clX}
\toprule
\# & Model & Özellik \\
\midrule
1 & Logistic Regression & Multinomial, StandardScaler, elasticNet CV \\
2 & Decision Tree       & maxDepth/minInstances CV, Gini impurity \\
3 & Random Forest       & 50--200 ağaç, bagging ensemble \\
4 & GBT (OneVsRest)     & N binary GBT, boosting ensemble \\
5 & Naive Bayes         & MinMaxScaler, multinomial/complement \\
\bottomrule
\end{tabularx}
\end{table}

Tüm modeller \texttt{CrossValidator} ile F1-Score optimizasyonu yapar.

\subsection{Değerlendirme Metrikleri}
Accuracy, Weighted F1, Precision, Recall, Log Loss ve $N \times N$ Confusion Matrix.

%=========================================================================
\section{MLflow Deney Takibi}

Her eğitim sonrası MLflow'a loglanan bilgiler:
\begin{itemize}[noitemsep]
  \item \textbf{Metrikler:} accuracy, f1, precision, recall, per-class P/R/F1
  \item \textbf{Parametreler:} tüm hiperparametreler, CV config
  \item \textbf{Artifact'lar:} confusion\_matrix.csv, feature importance (PNG+CSV), Spark model
  \item \textbf{Tag'ler:} classification\_type, num\_classes, model tipi
\end{itemize}

Backend: SQLite + artifact store. UI: \texttt{http://localhost:5000}.

%=========================================================================
\section{Streamlit Dashboard}

6 sayfalık interaktif web arayüzü:

\begin{table}[H]\centering
\begin{tabular}{cl}
\toprule
\# & Sayfa \\
\midrule
1 & Genel Bakış (sistem durumu) \\
2 & Veri Akışı (Kafka $\to$ Delta izleme) \\
3 & EDA (dağılım, korelasyon) \\
4 & Feature Engineering (5 feature detayı) \\
5 & Model Karşılaştırma (F1 sıralı) \\
6 & En İyi Model (production tavsiyesi) \\
\bottomrule
\end{tabular}
\end{table}

%=========================================================================
\section{Sonuç}

Proje, IoT ağ güvenliği için gerçek zamanlı veri toplama, temizleme, feature üretimi ve çok sınıflı sınıflandırmayı tek bir platform altında birleştirmiştir. 5 farklı model eğitilmiş, MLflow ile karşılaştırılmış ve Streamlit üzerinden görselleştirilmiştir. Gelecek adımlar: deep learning modelleri, real-time inference, otomatik alert sistemi.

%=========================================================================
\section*{Kaynakça}
\addcontentsline{toc}{section}{Kaynakça}
\begin{enumerate}[label={[\arabic*]},noitemsep]
  \item Ferrag, M.A. et al. ``Edge-IIoTset: A New Comprehensive Realistic Cyber Security Dataset.'' \textit{IEEE Access}, 2022.
  \item Zaharia, M. et al. ``Apache Spark: A Unified Engine for Big Data Processing.'' \textit{CACM}, 2016.
  \item Armbrust, M. et al. ``Delta Lake: High-Performance ACID Table Storage.'' \textit{VLDB}, 2020.
  \item Apache Kafka Documentation. \url{https://kafka.apache.org/documentation/}
  \item MLflow Documentation. \url{https://mlflow.org/docs/latest/index.html}
\end{enumerate}

\end{document}
